\documentclass{svproc}

\usepackage{url}
\def\UrlFont{\rmfamily}

\usepackage{graphicx}
\usepackage{amsmath}
\usepackage{array}

\sloppy

\newcommand{\maybegraphics}[2][]{%
  \IfFileExists{#2}{%
    \includegraphics[#1]{#2}%
  }{%
    \fbox{%
      \begin{minipage}{0.88\linewidth}
      \centering
      Missing figure file:\\
      \texttt{\detokenize{#2}}
      \end{minipage}
    }%
  }%
}

\begin{document}
\mainmatter

\title{Specimen-Leakage-Aware Evaluation for Shrimp Disease Instance Segmentation}

\titlerunning{Leakage-Aware Shrimp Disease Segmentation}

\author{
Vinh Dinh Nguyen\inst{1}
\and
Phong Van Nguyen\inst{1}
\and
Nhan Huu Tran\inst{1}
\and
Nhu Huynh Le Thi\inst{1}
}

\authorrunning{Vinh Dinh Nguyen et al.}

\tocauthor{Vinh Dinh Nguyen, Phong Van Nguyen, Nhan Huu Tran and Nhu Huynh Le Thi}

\institute{
FPT University Can Tho, Can Tho, Vietnam\\
\email{vinhnd29@fe.edu.vn, phongnv.ce192081@gmail.com, nhanthce181655@fpt.edu.vn, nhuhuynhlethi1508@gmail.com}
}

\maketitle

\begin{abstract}
Shrimp disease segmentation can support aquaculture monitoring by localizing visible disease regions rather than assigning only image-level labels. However, public shrimp disease datasets may contain multiple views of the same physical specimen. Random image-level splitting can therefore leak specimen-specific information across train and test sets and inflate evaluation results. This paper studies this problem for shrimp disease instance segmentation. We construct a hand-labeled instance segmentation dataset from ShrimpDiseaseImageBD, define specimen groups from the filename convention, and evaluate YOLO-based segmentation models under a disease-stratified grouped protocol. Across three seeds, random image-level splits create substantial train/test specimen overlap and higher mAP values, whereas grouped splitting eliminates overlap and gives a stricter estimate of unseen-specimen performance. Under the leakage-free protocol, YOLO11n-seg gives the best overall trade-off among the tested YOLO variants.
\keywords{shrimp disease, instance segmentation, data leakage, grouped split, YOLO}
\end{abstract}

\section{Introduction}

Computer vision is increasingly used for aquaculture health monitoring, where early recognition of visible disease symptoms can support faster intervention and reduce economic loss. In shrimp disease analysis, image classification can determine whether a disease is present, but it does not identify where disease evidence appears. Instance segmentation provides a more informative output by localizing visible disease regions, especially when symptoms are small, sparse, or mixed across disease types.

A major evaluation risk in shrimp disease images is repeated-specimen leakage. Public datasets may contain several images of the same physical shrimp captured from different views. If ordinary image-level random splitting is used, one view of a specimen can appear in training while another view appears in validation or testing. The model may then benefit from specimen-specific appearance, such as body texture, pose, lighting, scale, or background, rather than learning disease cues that generalize to unseen shrimp.

This paper addresses specimen leakage in shrimp disease instance segmentation. We use ShrimpDiseaseImageBD \cite{islam2025shrimpdiseaseimagebd} as the image source, construct a hand-labeled instance segmentation dataset for BG and WSSV disease regions, define specimen identities from the filename convention, and compare leakage-prone image-level splits with leakage-free grouped-specimen splits. The main contributions are: (i) a specimen-aware BG/WSSV instance segmentation dataset, (ii) an auditable grouped-specimen evaluation protocol, (iii) a three-seed analysis showing that image-level splitting inflates segmentation performance, and (iv) YOLO segmentation baselines with healthy-aware diagnostics.

\section{Related Work}

\subsection{Shrimp Disease Recognition and Localization}

ShrimpDiseaseImageBD is a public image dataset for computer-vision-based shrimp disease detection in Bangladesh \cite{islam2025shrimpdiseaseimagebd}. It contains four image-level classes: healthy shrimp, black gill (BG), white spot syndrome virus (WSSV), and \texttt{BG\_WSSV} co-infection. Recent shrimp-vision studies have explored transfer learning, augmentation, adversarial robustness, explainability, and lightweight architectures for disease recognition \cite{jone2026shrimpxnet,fei2025lightweight}. Most reported results focus on classification or detection. In contrast, this work focuses on instance-level disease localization and on whether the evaluation split is independent at the physical-specimen level.

\subsection{Leakage and Grouped Evaluation in Visual Learning}

Data leakage occurs when information unavailable at deployment time is indirectly used during model development or evaluation \cite{kaufman2012leakage}. In visual learning, leakage can arise from near-duplicate images, repeated subjects, repeated acquisition conditions, or repeated specimens across dataset splits. Similar concerns have been reported in medical imaging, where acquisition-site, patient, and repeated-scan confounding can inflate apparent performance \cite{zech2018confounding,roberts2021common,tampu2022inflation}. Shrimp disease imagery has an analogous specimen-level risk: multiple images can show the same shrimp, and all such images must remain in the same split to estimate unseen-specimen generalization.

\subsection{Lightweight YOLO-Based Segmentation}

YOLO-family models are widely used for real-time detection and segmentation because they provide a practical balance between accuracy, speed, and ease of training on custom datasets. This study evaluates YOLOv8, YOLO11, and YOLO26 segmentation variants using the Ultralytics framework \cite{ultralytics2026}. We focus on nano and small variants because shrimp disease segmentation should remain practical for iterative experimentation and future deployment on constrained hardware.

\section{Dataset and Leakage-Aware Evaluation Protocol}

\subsection{Source Dataset and Instance-Mask Annotation}

The image source is ShrimpDiseaseImageBD version 3 \cite{islam2025shrimpdiseaseimagebd}. The dataset contains 1149 raw images from 416 shrimp specimens, organized into four classes: 403 healthy, 198 BG, 328 WSSV, and 220 \texttt{BG\_WSSV} images. Each specimen has one to three images, which creates a risk of repeated-specimen leakage if images are split independently.

We relabeled the dataset for instance segmentation using two disease-region classes: BG and WSSV. Healthy images contain empty label files, while co-infection images can contain both BG and WSSV masks. Annotation was performed by a single annotator to reduce variation in mask boundary style, granularity, and disease-region interpretation. For repeated images of the same specimen, disease regions were annotated according to anatomical location rather than image position.
To make the segmentation labels more consistent, masks were assigned only to visible disease evidence rather than to the entire shrimp body. BG masks correspond to visually darkened gill regions, while WSSV masks correspond to visible white-spot or white-lesion regions on the shrimp body. Ambiguous shadows, reflections, background marks, and weak texture changes were not labeled as disease unless they were visually consistent with the target disease pattern. This conservative policy reduces noisy mask boundaries and makes healthy-image false positives measurable through empty-label images.

\begin{figure}
\centering
\maybegraphics[width=0.72\linewidth]{hand-labeled.png}
\caption{Specimen-level annotation consistency across repeated shrimp views. Red masks denote WSSV regions and blue masks denote BG regions.}
\label{fig:hand-labeled}
\end{figure}

\begin{table}
\caption{Statistics of the hand-labeled shrimp disease segmentation dataset.}
\label{tab:dataset_stats}
\begin{center}
\begin{tabular}{lr}
\hline
Property & Count \\
\hline
Images & 1149 \\
Specimens & 416 \\
Labeled diseased images & 746 \\
Healthy or empty-label images & 403 \\
Mask instances & 1031 \\
BG mask instances & 462 \\
WSSV mask instances & 569 \\
Co-infection images & 220 \\
Maximum images per specimen & 3 \\
\hline
\end{tabular}
\end{center}
\end{table}

\subsection{Specimen Group Definition}

The original filename convention follows \texttt{Disease-ShrimpID-img-ImageNumber.jpg}. We define a specimen group as \texttt{Disease::ShrimpID}, which groups all available views of the same shrimp specimen. Let $G_i$ denote the specimen group of image $i$ and $S_i$ denote the assigned split. A leakage-free split must satisfy:
\begin{equation}
G_i = G_j \Rightarrow S_i = S_j .
\end{equation}
Thus, all images of the same physical specimen must be assigned to exactly one of the training, validation, or test partitions.

\subsection{Split Policies and Leakage Audit}

We compare three split policies. The first is a plain random image split, where images are shuffled without disease stratification or specimen grouping. The second is a stratified random image split, where the class distribution is preserved across the four ShrimpDiseaseBD classes but the same specimen can still appear across splits. The third is a stratified grouped-specimen split, where specimens are grouped before splitting and disease-status stratification is applied at the group level.

Figure~\ref{fig:split_strategy} summarizes these policies. Image-level splitting assigns repeated views independently and can place the same specimen in multiple partitions. The grouped-specimen protocol first collects all views of a physical specimen into one group and then assigns the full group to training, validation, or testing.

\begin{figure}
\centering
\maybegraphics[width=0.96\linewidth]{data_splitting_strategies_using_specimen.png}
\caption{Data-splitting strategies using specimen groups. Plain random and stratified random image-level splits can leak repeated views of the same shrimp across partitions. Stratified grouped-specimen splitting assigns all views of each specimen to one partition while preserving class balance at the group level.}
\label{fig:split_strategy}
\end{figure}

The grouped-specimen split is the main evaluation protocol. The two image-level splits are retained only to quantify leakage-prone evaluation. For every run, we save a split manifest, split summary, and split fingerprint, and compute train/test and validation/test specimen overlap. A grouped split is accepted only when both overlap values are zero. 
This audit is important because class-balanced splitting alone does not guarantee specimen independence. A split can preserve the four disease-class proportions while still placing repeated views of the same shrimp in both training and testing. Therefore, the split unit must match the biological unit of generalization: the physical specimen rather than the individual image. In our experiments, the image-level splits are used only to quantify leakage-prone behavior, while the grouped split is treated as the main protocol for final model comparison.

\subsection{Evaluation Metrics}

We report metrics on the full test set and on the labeled-only test set. The full test includes both diseased and healthy images, while the labeled-only test includes only images with disease masks. The primary metrics are mask mAP50 and mask mAP50--95; box mAP is recorded for completeness. We additionally report healthy mask false-positive rate, healthy false-positive masks per image, disease missed-image rate, and mask count mean absolute error (MAE).

For validation-based model selection, we use a healthy-aware score:
\begin{equation}
S =
\mathrm{mAP50}_{\mathrm{lab}}
- 0.05 \cdot \mathrm{MAE}_{\mathrm{count}}
- 0.15 \cdot R_{\mathrm{miss}}
- 0.10 \cdot R_{\mathrm{healthyFP}},
\end{equation}
where $\mathrm{mAP50}_{\mathrm{lab}}$ is labeled-only mask mAP50, $R_{\mathrm{miss}}$ is the disease missed-image rate, and $R_{\mathrm{healthyFP}}$ is the healthy false-positive rate. Test-set healthy-aware scores are reported only for analysis.

\section{Experimental Setup}

\subsection{YOLO Segmentation Baselines}

We evaluate YOLOv8-seg, YOLO11-seg, and YOLO26-seg nano/small variants. The model sweep is conducted under the stratified grouped-specimen protocol to avoid selecting a backbone using leakage-prone behavior. After the sweep, the best trade-off model is fixed for preprocessing and augmentation ablations.

\subsection{Training Configuration}

All models are trained using the Ultralytics segmentation framework \cite{ultralytics2026}. Images are resized to 640 pixels. Models are trained for up to 100 epochs with patience 30. The default baseline uses a clean-light augmentation policy: horizontal flip probability 0.5, HSV perturbation $(h=0.01,s=0.35,v=0.20)$, translation 0.05, scale 0.20, and no mosaic, mixup, cutmix, or copy-paste. The best checkpoint is selected using validation metrics, and test metrics are computed only after checkpoint selection.

\subsection{Preprocessing and Augmentation Ablations}

We evaluate copy-paste, train-only photometric augmentation, noise/blur/compression augmentation, Fourier high-pass preprocessing, and wavelet-based enhancement. Fourier and wavelet variants are tested either as train-only transformed copies or as deterministic preprocessing applied to all splits. These experiments are treated as ablations rather than assumed improvements.

\section{Results}

\subsection{Image-Level Splits Inflate Segmentation Performance}

Table~\ref{tab:split_policy_summary} summarizes YOLO11n-seg across three seeds under the three split policies. Both image-level policies produce substantial train/test specimen overlap. Plain random splitting creates $98.7 \pm 1.2$ overlapping specimen groups on average, and stratified random splitting creates $94.0 \pm 6.1$. In contrast, the stratified grouped-specimen split keeps train/test overlap at zero for every seed.

The leakage-prone image-level splits report higher mAP values than the grouped split. Plain random splitting reaches $0.677 \pm 0.013$ labeled-only mask mAP50, while the grouped split reaches $0.473 \pm 0.035$. This gap indicates that image-level test results are inflated by repeated-specimen overlap and do not estimate unseen-specimen generalization.

\begin{table}
\caption{Three-seed split-policy summary using YOLO11n-seg. Values are mean $\pm$ standard deviation.}
\label{tab:split_policy_summary}
\begin{center}
\scriptsize
\setlength{\tabcolsep}{3pt}
\renewcommand{\arraystretch}{1.08}
\resizebox{\linewidth}{!}{%
\begin{tabular}{lrrrrrr}
\hline
Policy & Leak & Full50 & Lab50 & Lab50--95 & HFP & HScore \\
\hline
Plain image & $98.7{\pm}1.2$ & $0.671{\pm}0.013$ & $0.677{\pm}0.013$ & $0.299{\pm}0.024$ & $0.054{\pm}0.027$ & $0.648{\pm}0.012$ \\
Stratified image & $94.0{\pm}6.1$ & $0.651{\pm}0.060$ & $0.659{\pm}0.062$ & $0.282{\pm}0.016$ & $0.098{\pm}0.024$ & $0.628{\pm}0.053$ \\
Grouped specimen & $0.0{\pm}0.0$ & $0.451{\pm}0.027$ & $0.473{\pm}0.035$ & $0.144{\pm}0.021$ & $0.236{\pm}0.110$ & $0.415{\pm}0.039$ \\
\hline
\end{tabular}%
}
\end{center}
\end{table}

Stratified random splitting has a lower mean score and higher variance than plain random splitting, but it still contains specimen leakage. This shows that class-distribution balancing alone is insufficient. For shrimp disease segmentation, split independence must be enforced at the physical-specimen level.

\subsection{Leakage-Free YOLO Baseline Comparison}

Table~\ref{tab:model_sweep} reports the YOLO model sweep under the stratified grouped-specimen split. YOLO11n-seg obtains the best full-test mask mAP50, labeled-only mask mAP50, and healthy-aware score among the tested variants. YOLO26s-seg gives the highest mask mAP50--95, and YOLO26n-seg has the lowest healthy false-positive rate, but neither provides the strongest overall trade-off.

\begin{table}
\caption{YOLO model sweep under the stratified grouped-specimen split.}
\label{tab:model_sweep}
\begin{center}
\scriptsize
\setlength{\tabcolsep}{4pt}
\renewcommand{\arraystretch}{1.08}
\resizebox{\linewidth}{!}{%
\begin{tabular}{lrrrrrr}
\hline
Model & Params (M) & Full50 & Lab50 & Mask50--95 & HFP & HScore \\
\hline
YOLO11n-seg & 2.877 & 0.482 & 0.512 & 0.168 & 0.244 & 0.460 \\
YOLO26s-seg & 11.506 & 0.417 & 0.460 & 0.182 & 0.268 & 0.397 \\
YOLOv8n-seg & 3.410 & 0.395 & 0.432 & 0.139 & 0.341 & 0.353 \\
YOLOv8s-seg & 11.821 & 0.391 & 0.444 & 0.159 & 0.341 & 0.364 \\
YOLO26n-seg & 3.126 & 0.430 & 0.463 & 0.160 & 0.195 & 0.393 \\
YOLO11s-seg & 10.113 & 0.297 & 0.361 & 0.124 & 0.512 & 0.250 \\
\hline
\end{tabular}%
}
\end{center}
\end{table}

Although YOLO11n-seg is the best tested baseline, the absolute grouped-split scores remain moderate. The main result is therefore not that YOLO solves shrimp disease segmentation, but that leakage-free evaluation gives a stricter and more realistic estimate of performance.

\subsection{Effect of Preprocessing and Augmentation}

Table~\ref{tab:preproc_ablation} summarizes representative preprocessing and augmentation ablations using YOLO11n-seg under the grouped split. None of the tested variants improves the healthy-aware baseline. Copy-paste slightly increases labeled-only mAP50 from 0.512 to 0.515, but it doubles the healthy false-positive rate from 0.244 to 0.488 and lowers the healthy-aware score.

\begin{table}
\caption{Preprocessing and augmentation ablations using YOLO11n-seg under the grouped split.}
\label{tab:preproc_ablation}
\begin{center}
\scriptsize
\setlength{\tabcolsep}{3pt}
\renewcommand{\arraystretch}{1.08}
\resizebox{\linewidth}{!}{%
\begin{tabular}{lrrrrr}
\hline
Method & Full50 & Lab50 & Mask50--95 & HFP & HScore \\
\hline
Clean-light baseline & 0.482 & 0.512 & 0.168 & 0.244 & 0.460 \\
Copy-paste 0.10 mixup & 0.464 & 0.515 & 0.144 & 0.488 & 0.424 \\
Photometric medium & 0.443 & 0.477 & 0.168 & 0.244 & 0.420 \\
Noise/blur/compression & 0.449 & 0.470 & 0.136 & 0.366 & 0.401 \\
Fourier, train only & 0.436 & 0.468 & 0.163 & 0.341 & 0.403 \\
Fourier, all splits & 0.452 & 0.498 & 0.166 & 0.244 & 0.426 \\
Wavelet db2, all splits & 0.406 & 0.435 & 0.131 & 0.293 & 0.352 \\
\hline
\end{tabular}%
}
\end{center}
\end{table}

These results suggest that simple image enhancement is not automatically beneficial for shrimp disease segmentation. Frequency-domain preprocessing may sharpen local texture, but it can also alter global appearance and increase false alarms. Preprocessing should therefore be selected using both disease-localization metrics and healthy false-positive diagnostics.

\section{Discussion and Limitations}

The experiments support three observations. First, specimen leakage is not a minor implementation detail. Across three seeds, image-level random splits place many physical specimens across train and test partitions. The resulting metrics are higher, but they do not reflect unseen-specimen generalization. Grouped splitting lowers the reported scores because it removes this shortcut.

Second, stratified grouped splitting is fairer but more difficult on small datasets. Once all views of a shrimp must move together, fewer balanced train/validation/test combinations are possible. Test difficulty can vary depending on which specimens are assigned to the test set. For this reason, seed summaries are more informative than a single split result.

Third, healthy images must be evaluated explicitly. A segmentation model can obtain reasonable disease mAP while still producing false disease masks on healthy shrimp. In a farm-monitoring setting, such false positives can create unnecessary inspection or treatment costs. The healthy-aware score is therefore used as a diagnostic summary, not as a replacement for standard mAP.

This study has several limitations. The dataset is derived from a single public source and does not provide external validation across farms, cameras, or regions. The masks were produced by a single annotator, which improves internal consistency but does not measure inter-annotator uncertainty. The grouped-split results show that absolute segmentation performance remains moderate, especially for mask mAP50--95. Finally, the current experiments evaluate offline images only; real farm deployment would require testing under uncontrolled lighting, water, background, camera, and mobile-device conditions.

The labeled instance segmentation dataset is available as a Roboflow Universe dataset \cite{nguyen2026shrimpseg}, and Roboflow was used for annotation management and YOLO-format export \cite{roboflow}.

\section{Conclusion}

This paper presents a specimen-leakage-aware evaluation study for shrimp disease instance segmentation. We construct a hand-labeled BG/WSSV segmentation dataset from ShrimpDiseaseImageBD, define specimen groups from the filename convention, and enforce a disease-stratified grouped split to prevent same-shrimp train/test leakage. The split-policy experiments show that image-level splitting creates substantial specimen overlap and inflates apparent segmentation performance. Under the leakage-free protocol, YOLO11n-seg provides the best overall baseline among the tested YOLO variants, but the grouped-split results also show that shrimp disease segmentation remains challenging. Future work should expand cross-farm validation, quantify annotation uncertainty, improve robustness under real imaging conditions, and evaluate deployment-oriented models on mobile or edge hardware.

\bibliographystyle{splncs03_unsrt}
\bibliography{references}

\end{document}

